% =========================================================
% Zorn's Lemma and the Axiom of Choice
% Source: Mark Dean, Order Theory (Brown University, 2015)
% =========================================================

\subsection{Zorn's Lemma and the Axiom of Choice}

% ---------------------------------------------------------
% TOOLKIT
% ---------------------------------------------------------

\begin{remark*}[Why AC is non-constructive]
For finite collections, the axiom is provable without extra assumptions
(see SRF-STO-C08-S8.1). For infinite collections, no construction can
produce the choice function: AC asserts its existence without exhibiting it.
This non-constructive character is both AC's power and the source of the
foundational controversy surrounding it.
\end{remark*}

% ---------------------------------------------------------
% Equivalent formulations
% ---------------------------------------------------------
\begin{theorembox}{Theorem (Alternative Forms of the Axiom of Choice)}
\label{thm:ac-alternative-forms}
Over ZF set theory, the following statements are equivalent.
\begin{enumerate}
  \item \textbf{Axiom of Choice.} Every indexed family of nonempty sets has a
        choice function.
  \item \textbf{Pairwise-disjoint choice form.} If
        $\{A_\lambda\}_{\lambda\in\Lambda}$ is an indexed family of nonempty
        pairwise disjoint sets, then there is a set
        $B\subseteq\bigcup_{\lambda\in\Lambda}A_\lambda$ such that
        $B\cap A_\lambda$ has exactly one element for each
        $\lambda\in\Lambda$.
  \item \textbf{Zorn's Lemma.} If each chain in a nonempty partially ordered
        set has an upper bound, then the ordered set has a maximal element.
  \item \textbf{Zermelo's Well-Ordering Theorem.} Every set can be
        well-ordered.
  \item \textbf{Tukey's Lemma.} Every nonempty family $\mathcal{F}$ of sets of
        finite character has a maximal element, where finite character means
        that $X\in\mathcal{F}$ exactly when every finite subset of $X$ belongs
        to $\mathcal{F}$.
\end{enumerate}
\end{theorembox}

\begin{remark*}[How to read the equivalences]
These are not different axioms for this book; they are interchangeable tools.
The usual selection-function statement is the official axiom. Zorn's Lemma is
often the most useful proof engine, the well-ordering theorem is the order
theoretic form, and Tukey's Lemma packages maximality arguments for families
controlled by their finite subfamilies.
\end{remark*}

% ---------------------------------------------------------
% Zorn's Lemma
% ---------------------------------------------------------
\begin{theorembox}{Zorn's Lemma}
\label{thm:zorn}
Let $(P, \leq)$ be a poset in which every chain (loset) has an upper bound
in $P$. Then $P$ has at least one maximal element.
\end{theorembox}

\begin{remark*}[Reading Zorn's Lemma]
The hypothesis ``every chain has an upper bound'' does \emph{not} require
the bound to lie in the chain itself. The conclusion gives a maximal
element: an element above which nothing in $P$ sits. In a partial order,
there may be many maximal elements; Zorn guarantees at least one.
\end{remark*}

\begin{remark*}[Typical proof pattern using Zorn's Lemma]
To show an object of type $X$ exists with a maximal property:
\begin{enumerate}
  \item Partially order candidate objects by inclusion (or extension).
  \item Verify every chain of candidates has an upper bound (usually its union).
  \item Conclude by Zorn that a maximal candidate exists.
  \item Show the maximal candidate has the desired property (often: if it
        lacked it, it could be enlarged, contradicting maximality).
\end{enumerate}
This pattern recurs throughout algebra and analysis: maximal ideals,
algebraic bases (Hamel bases), maximal consistent sets, and Szpilrajn's Theorem.
\end{remark*}

% ---------------------------------------------------------
% Hausdorff Maximum Principle
% ---------------------------------------------------------
\begin{theorembox}{Hausdorff Maximum Principle}
\label{thm:hausdorff}
In every poset $(P, \leq)$ there exists a $\subseteq$-maximal chain; that
is, a chain $C \subseteq P$ such that no chain $C' \subseteq P$ strictly
contains $C$.
\end{theorembox}

\begin{remark*}[Relation to Zorn]
The Hausdorff Maximum Principle is an immediate corollary of Zorn's Lemma:
take the poset of chains ordered by inclusion; every chain of chains is
bounded by its union; Zorn gives a maximal chain. Conversely, Hausdorff
implies Zorn. Both are equivalent to AC over ZF.
\end{remark*}

\begin{remark*}[Application: Szpilrajn via Hausdorff]
In the proof of Szpilrajn's Theorem (\hyperref[thm:szpilrajn]{$\uparrow$ Thm}), let
$T_X$ be the set of all partial orders on $X$ extending $\succsim$, ordered
by $\subseteq$. By the Hausdorff Maximum Principle, $T_X$ has a maximal chain
$A$; its union $\succsim^* = \bigcup A$ is a partial order extending
$\succsim$. If $\succsim^*$ were not total, there would exist incomparable
$x, y$ and one could extend $\succsim^*$ by adding $(x,y)$ (closing under
transitivity), contradicting maximality of $A$. Hence $\succsim^*$ is a
linear order. \qed
\end{remark*}

\begin{remark*}[Transition]
Zorn's Lemma is the workhorse for existence proofs throughout abstract
algebra and analysis. It first appears in the project context here and at
the proof stubs in \S11 of the Chapter~VIII exercises.
\end{remark*}
